% GENERATED FILE. Edit canonical lessons, then run npm run generate:books.
% canonical-source-hash: fnv1a64:976e999ec944eecc
% canonical-lessons: ES-C683-region, ES-C683-comarca, ES-C683-municipio, ES-C683-aldea, ES-C683-alcaldia

\chapter{Regiones y pueblos}
\label{ch:es-regiones-y-pueblos}
\hlchaptermodality{683}

\begin{chapteropening}
\textbf{By the end of this chapter:} \emph{I can say which region, district or municipality a place is in, find the mayor’s office and talk about hamlets.}

Complete the last lesson of chapter 683: I can say which region, district or municipality a place is in, find the mayor’s office and talk about hamlets.
\end{chapteropening}

\section[la región]{la región — the region}

\label{lesson:ES-C683-region}



\begin{quote}
Before the new one: say the Spanish for the kerb, then the Spanish for pedestrian.
\end{quote}

\subsection*{You'll want to know: la región}
\textbf{la región} — \textquotedblleft{}the region\textquotedblright{}.

\begin{quote}
Es una región famosa por su vino.
\end{quote}

\begin{quote}
\emph{It is a region famous for its wine.}
\end{quote}

\emph{La región} is wider than \emph{la zona}.

\begin{cousinweb}
\emph{región} is a learned borrowing from Latin \emph{regiō} \textquotedblleft{}a direction, a boundary, a district\textquotedblright{} (EtymDB), from \emph{regere} \textquotedblleft{}to guide, to rule\textquotedblright{}.
\end{cousinweb}

\begin{grammarlens}[title={What you've built}]
One more word for telling, explaining and describing at B1.
\end{grammarlens}

\subsection*{Your turn}
\begin{itemize}
\raggedright
  \item \emph{Say it:} \emph{la región}
  \item \emph{Say it:} \emph{la región}, once more
  \item \emph{Say it:} a sentence of your own with \emph{la región}
  \item \emph{Recall:} say the Spanish for the kerb, then the Spanish for pedestrian, then say \emph{la región} again
\end{itemize}

\begin{tcolorbox}[breakable,colback=teal!4,colframe=teal!35!black,title={Before you move on}]
What does la región mean? (\textquotedblleft{}The region\textquotedblright{}.) Say it once more.
\end{tcolorbox}
\section[la comarca]{la comarca — the district, the area (of several villages)}

\label{lesson:ES-C683-comarca}



\begin{quote}
Before the new one: say the Spanish for pedestrian, then the Spanish for the region.
\end{quote}

\subsection*{You'll want to know: la comarca}
\textbf{la comarca} — \textquotedblleft{}the district, the area (of several villages)\textquotedblright{}.

\begin{quote}
Es el pueblo más grande de la comarca.
\end{quote}

\begin{quote}
\emph{It is the biggest village in the district.}
\end{quote}

Smaller than a province, bigger than a municipality.

\begin{cousinweb}
\emph{comarca} is usually traced to \emph{marca}, an old word for a border district, which is of Germanic origin.
\end{cousinweb}

\begin{grammarlens}[title={What you've built}]
One more word for telling, explaining and describing at B1.
\end{grammarlens}

\subsection*{Your turn}
\begin{itemize}
\raggedright
  \item \emph{Say it:} \emph{la comarca}
  \item \emph{Say it:} \emph{la comarca}, once more
  \item \emph{Say it:} a sentence of your own with \emph{la comarca}
  \item \emph{Recall:} say the Spanish for pedestrian, then the Spanish for the region, then say \emph{la comarca} again
\end{itemize}

\begin{tcolorbox}[breakable,colback=teal!4,colframe=teal!35!black,title={Before you move on}]
What does la comarca mean? (\textquotedblleft{}The district, the area (of several villages)\textquotedblright{}.) Say it once more.
\end{tcolorbox}
\section[el municipio]{el municipio — the municipality, the town council}

\label{lesson:ES-C683-municipio}



\begin{quote}
Before the new one: say the Spanish for the region, then the Spanish for the district.
\end{quote}

\subsection*{You'll want to know: el municipio}
\textbf{el municipio} — \textquotedblleft{}the municipality, the town council\textquotedblright{}.

\begin{quote}
El municipio tiene tres pueblos.
\end{quote}

\begin{quote}
\emph{The municipality has three villages.}
\end{quote}

Its people and its council: \emph{el ayuntamiento} runs \emph{el municipio}.

\begin{cousinweb}
\emph{municipio} is a learned borrowing from Latin \emph{mūnicipium} \textquotedblleft{}a free town with its own rights\textquotedblright{} (EtymDB).
\end{cousinweb}

\begin{grammarlens}[title={What you've built}]
One more word for telling, explaining and describing at B1.
\end{grammarlens}

\subsection*{Your turn}
\begin{itemize}
\raggedright
  \item \emph{Say it:} \emph{el municipio}
  \item \emph{Say it:} \emph{el municipio}, once more
  \item \emph{Say it:} a sentence of your own with \emph{el municipio}
  \item \emph{Recall:} say the Spanish for the region, then the Spanish for the district, then say \emph{el municipio} again
\end{itemize}

\begin{tcolorbox}[breakable,colback=teal!4,colframe=teal!35!black,title={Before you move on}]
What does el municipio mean? (\textquotedblleft{}The municipality, the town council\textquotedblright{}.) Say it once more.
\end{tcolorbox}
\section[la aldea]{la aldea — the hamlet, the small village}

\label{lesson:ES-C683-aldea}



\begin{quote}
Before the new one: say the Spanish for the district, then the Spanish for the municipality.
\end{quote}

\subsection*{You'll want to know: la aldea}
\textbf{la aldea} — \textquotedblleft{}the hamlet, the small village\textquotedblright{}.

\begin{quote}
En la aldea viven veinte personas.
\end{quote}

\begin{quote}
\emph{Twenty people live in the hamlet.}
\end{quote}

Smaller than \emph{un pueblo}.

\begin{cousinweb}
\emph{aldea} comes from Arabic \emph{aḍ-ḍayʿa} \textquotedblleft{}the farm, the village\textquotedblright{} (EtymDB), with the Arabic article \emph{al-} fused on.
\end{cousinweb}

\begin{grammarlens}[title={What you've built}]
One more word for telling, explaining and describing at B1.
\end{grammarlens}

\subsection*{Your turn}
\begin{itemize}
\raggedright
  \item \emph{Say it:} \emph{la aldea}
  \item \emph{Say it:} \emph{la aldea}, once more
  \item \emph{Say it:} a sentence of your own with \emph{la aldea}
  \item \emph{Recall:} say the Spanish for the district, then the Spanish for the municipality, then say \emph{la aldea} again
\end{itemize}

\begin{tcolorbox}[breakable,colback=teal!4,colframe=teal!35!black,title={Before you move on}]
What does la aldea mean? (\textquotedblleft{}The hamlet, the small village\textquotedblright{}.) Say it once more.
\end{tcolorbox}
\section[la alcaldía]{la alcaldía — the mayor’s office; the office of mayor}

\label{lesson:ES-C683-alcaldia}



\begin{quote}
Before the new one: say the Spanish for the municipality, then the Spanish for the hamlet.
\end{quote}

\subsection*{You'll want to know: la alcaldía}
\textbf{la alcaldía} — \textquotedblleft{}the mayor’s office; the office of mayor\textquotedblright{}.

\begin{quote}
La alcaldía está en la plaza mayor.
\end{quote}

\begin{quote}
\emph{The mayor’s office is on the main square.}
\end{quote}

Built on \emph{el alcalde}, the mayor.

\begin{cousinweb}
\emph{alcaldía} is built on \emph{alcalde} \textquotedblleft{}mayor\textquotedblright{}, which comes from Arabic \emph{al-qāḍī} \textquotedblleft{}the judge\textquotedblright{} (EtymDB).
\end{cousinweb}

\begin{grammarlens}[title={What you've built}]
One more word for telling, explaining and describing at B1.
\end{grammarlens}

\subsection*{Your turn}
\begin{itemize}
\raggedright
  \item \emph{Say it:} \emph{la alcaldía}
  \item \emph{Say it:} \emph{la alcaldía}, once more
  \item \emph{Say it:} a sentence of your own with \emph{la alcaldía}
  \item \emph{Recall:} say the Spanish for the municipality, then the Spanish for the hamlet, then say \emph{la alcaldía} again
\end{itemize}

\begin{tcolorbox}[breakable,colback=teal!4,colframe=teal!35!black,title={Before you move on}]
What does la alcaldía mean? (\textquotedblleft{}The mayor’s office; the office of mayor\textquotedblright{}.) Say it once more.
\end{tcolorbox}
